\ifdefined\gkver\else\def\gkver{student}\fi
\documentclass[\gkver]{gaokaozhenti}
\begin{document}

\chapter{2023年上海}

\begin{enumerate}
\item
%% number: 1
%% paperName: 2023年上海高考第 1 题 3 分
%% typeId: 1
%% type: 单选题
%% chapter: 原子和原子核
%% point: 原子核式结构
%% method: 无
%% score: 3
%% degree: 445
%% duplicateId: 0
%% body:
$\alpha$ 粒子散射实验\xzanswer{A}
\fourchoices[answer=A]
{需要在真空中完成}
{使用荧光屏是为了阻挡 $\alpha$ 粒子}
{所用显微镜需始终正对放射源}
{证明了原子核中存在质子}
%% answer: A
%% memo:
\memoanswer{
本题原卷及材料中未提供详解，待补充。
}

\item
%% number: 2
%% paperName: 2023年上海高考第 2 题 3 分
%% typeId: 1
%% type: 单选题
%% chapter: 电路
%% point: 串并联电路
%% method: 无
%% score: 3
%% degree: 870
%% duplicateId: 0
%% body:
将四个完全相同的灯泡按如图所示的电路连接，闭合开关后各灯泡均发光，则\xzanswer{A}
\onepicture[w=5.5cm, align=right]{\includesvg{figs/02}}
\fourchoices[answer=A]
{灯泡 $L_{1}$ 最亮}
{灯泡 $L_{2}$ 最亮}
{灯泡 $L_{3}$ 最亮}
{四个灯泡一样亮}
%% answer: A
%% memo:
\memoanswer{
本题原卷及材料中未提供详解，待补充。
}

\item
%% number: 3
%% paperName: 2023年上海高考第 3 题 3 分
%% typeId: 1
%% type: 单选题
%% chapter: 功和能
%% point: 机械能守恒定律
%% method: 无
%% score: 3
%% degree: 710
%% duplicateId: 0
%% body:
一个爆竹在空中炸成多块，其中 a、b、c 三块的质量分别为 $m$、$2m$、$3m$，它们在相同高度以相同的速率分别沿水平、向上、向下方向，空气阻力不计，三者落到水平地面时\xzanswer{D}
\fourchoices[answer=D]
{a 的速率最大}
{b 的速率最大}
{c 的速率最大}
{速率一样大}
%% answer: D
%% memo:
\memoanswer{
本题原卷及材料中未提供详解，待补充。
}

\item
%% number: 4
%% paperName: 2023年上海高考第 4 题 3 分
%% typeId: 1
%% type: 单选题
%% chapter: 气体性质
%% point: 气体图象
%% method: 无
%% score: 3
%% degree: 930
%% duplicateId: 0
%% body:
一定质量的理想气体压强 $p$ 随体积 $V$ 的变化关系如图所示，直线 ad、bc 均与 $p$ 轴平行，曲线 ab 和 cd 均为反比例函数曲线的一部分。气体处于 a、b、c、d 四个状态时的温度分别为 $T_{a}$、$T_{b}$、$T_{c}$ 和 $T_{d}$，则\xzanswer{B}
\onepicture[w=4.5cm, align=right]{\includesvg{figs/04}}
\fourchoices[answer=B]
{$T_{a} > T_{b}$}
{$T_{b} > T_{c}$}
{$T_{c} > T_{d}$}
{$T_{d} > T_{a}$}
%% answer: B
%% memo:
\memoanswer{
本题原卷及材料中未提供详解，待补充。
}

\item
%% number: 5
%% paperName: 2023年上海高考第 5 题 3 分
%% typeId: 1
%% type: 单选题
%% chapter: 直线运动
%% point: 平均速度和瞬时速度
%% method: 无
%% score: 3
%% degree: 830
%% duplicateId: 0
%% body:
某次 $100 \Um$ 短跑比赛，发令枪响后 $3.298 \Us$ 时各运动员的位置如图。位于第 3 道的运动员前 $30 \Um$ 的平均速度为 $v_{1}$，在图示时刻的速度为 $v_{2}$，由图中信息\xzanswer{B}
\onepicture[w=9cm, align=right]{\includesvg{figs/05}}
\fourchoices[answer=B]
{能得出 $v_{1}$ 和 $v_{2}$}
{只能得出 $v_{1}$}
{不能得出 $v_{1}$ 和 $v_{2}$}
{只能得出 $v_{2}$}
%% answer: B
%% memo:
\memoanswer{
数据来自于“苏炳添 2021 年东京奥运会百米半决赛”。上网查了下，苏炳添半决赛跑至 $30 \Um$ 处用时 $3.73 \Us$，是目前世界第一，去除起跑反应时间 $0.142 \Us$，也是 $3.58 \Us$，题目中的 $3.298 \Us$ 来源不明。

\onepicture[w=8cm, align=right]{figs/05_an.jpg}
}

\item
%% number: 6
%% paperName: 2023年上海高考第 6 题 3 分
%% typeId: 1
%% type: 单选题
%% chapter: 电场
%% point: 电荷守恒定律
%% method: 无
%% score: 3
%% degree: 920
%% duplicateId: 0
%% body:
真空中有完全相同的导体球 $X$、$Y$、$Z$，其中 $X$ 带电 $+4 \UuC$、$Y$ 不带电，$Z$ 带电 $-10 \UuC$。先将 $Y$ 与 $X$ 接触，分开后 $Y$ 再与 $Z$ 接触，最后 $Y$ 带电\xzanswer{A}
\fourchoices[answer=A]
{$-4 \UuC$}
{$-3 \UuC$}
{$-2 \UuC$}
{$+4 \UuC$}
%% answer: A
%% memo:
\memoanswer{
本题原卷及材料中未提供详解，待补充。
}

\item
%% number: 7
%% paperName: 2023年上海高考第 7 题 3 分
%% typeId: 1
%% type: 单选题
%% chapter: 机械波
%% point: 波长频率波速
%% method: 无
%% score: 3
%% degree: 840
%% duplicateId: 0
%% body:
一周期为 $T$、向右传播的简谐横波在 $t = 0 \Us$ 时传到 P 点。波形如图所示，在 $t = 8T$ 时，a、b 两点间的波形为\xzanswer{C}
\onepicture[w=9cm, align=right]{\includesvg{figs/07a}}
\fourchoices[answer=C, ispicture=true, h=1.6cm]
{figs/07b.png}
{figs/07c.png}
{figs/07d.png}
{figs/07e.png}
%% answer: C
%% memo:
\memoanswer{
本题原卷及材料中未提供详解，待补充。
}

\item
%% number: 8
%% paperName: 2023年上海高考第 8 题 3 分
%% typeId: 1
%% type: 单选题
%% chapter: 电场
%% point: 电场中的图象
%% method: 无
%% score: 3
%% degree: 790
%% duplicateId: 0
%% body:
用“光镊”技术可将纳米颗粒悬浮在真空中的某点 P。当该微粒偏离 P 很小距离 $x$ 时，其所受合力大小与 $x$ 成正比，方向始终指向 P。现使位于 P 的颗粒获得一微小速度，则该微粒在此后的运动过程中可能\xzanswer{B}
\fourchoices[answer=B]
{速度减小时加速度减小}
{速度增大时加速度减小}
{速度减小时加速度不变}
{速度增大时加速度不变}
%% answer: B
%% memo:
\memoanswer{
本题原卷及材料中未提供详解，待补充。
}

\item
%% number: 9
%% paperName: 2023年上海高考第 9 题 4 分
%% typeId: 1
%% type: 单选题
%% chapter: 电场
%% point: 电势能电势差电场力做功
%% method: 无
%% score: 4
%% degree: 730
%% duplicateId: 0
%% body:
某电场中 $x$ 轴上 $0 \leq x \leq a$ 区间电势 $\varphi$ 随 $x$ 变化的关系如图所示。在将一带负电的点电荷沿 $x$ 轴从 $x = 0$ 移动到 $x = a$ 的过程中，其电势能 $E_{p}$ 随 $x$ 变化的关系可能是\xzanswer{C}
\onepicture[w=4cm, align=right]{\includesvg{figs/09a}}
\fourchoices[answer=C, ispicture=true, h=2.6cm]
{figs/09b.png}
{figs/09c.png}
{figs/09d.png}
{figs/09e.png}
%% answer: C
%% memo:
\memoanswer{
本题原卷及材料中未提供详解，待补充。
}

\item
%% number: 10
%% paperName: 2023年上海高考第 10 题 4 分
%% typeId: 1
%% type: 单选题
%% chapter: 直线运动
%% point: 匀速直线运动
%% method: 近似与估算
%% score: 4
%% degree: 740
%% duplicateId: 0
%% body:
坦克炮管通常在发射数百发炮弹后需报废，炮弹离开炮管时的速度可达 $1\,000 \Ums$，一根炮管在报废前发射的所有炮弹在该炮管中运动的累积时间约为\xzanswer{A}
\fourchoices[answer=A]
{5 秒}
{5 分钟}
{5 小时}
{5 个月}
%% answer: A
%% memo:
\memoanswer{
本题原卷及材料中未提供详解，待补充。
}

\item
%% number: 11
%% paperName: 2023年上海高考第 11 题 4 分
%% typeId: 1
%% type: 单选题
%% chapter: 电场
%% point: 带电粒子的平衡
%% method: 无
%% score: 4
%% degree: 840
%% duplicateId: 0
%% body:
带电量分别为 $+q$、$+q$ 和 $-2q$ 的点电荷固定在等边三角形的三个顶点上，构成一带电体。将其先后置于如图（a）、（b）所示的匀强电场 $E_{a}$ 和非匀强电场 $E_{b}$ 中，该带电体受到的电场力大小分别为\xzanswer{D}
\twopicture[h=3.5cm, align=right]
{figs/11a.png}
{figs/11b.png}
\fourchoices[answer=D]
{$F_{a} = 0$，$F_{b} = 0$}
{$F_{a} \neq 0$，$F_{b} = 0$}
{$F_{a} \neq 0$，$F_{b} \neq 0$}
{$F_{a} = 0$，$F_{b} \neq 0$}
%% answer: D
%% memo:
\memoanswer{
本题原卷及材料中未提供详解，待补充。
}

\item
%% number: 12
%% paperName: 2023年上海高考第 12 题 4 分
%% typeId: 1
%% type: 单选题
%% chapter: 电磁感应
%% point: 电磁感应中的匀变速
%% method: 无
%% score: 4
%% degree: 610
%% duplicateId: 0
%% body:
如图（a），两根互相平行、电阻不计的光滑导轨位于水平面内，左端用电阻连接，处于竖直向下的匀强磁场中。一金属棒置于导轨上，与导轨垂直且接触良好，受到一垂直于棒的水平外力 $F$ 作用。以向右为正方向，棒所受的安培力 $F_{A}$ 与时间 $t$ 的关系如图（b）所示。则外力 $F$ 随时间 $t$ 的关系可能是\xzanswer{C}
\twopicture[h=3.5cm, align=right]
{figs/12a.png}
{figs/12b.png}
\fourchoices[answer=C, ispicture=true, h=2.6cm]
{figs/12c.png}
{figs/12d.png}
{figs/12e.png}
{figs/12f.png}
%% answer: C
%% memo:
\memoanswer{
由 $F_{A} = \frac{B^{2}L^{2}v}{R}$ 和图（b）可知，金属棒在 $0 \sim t_{0}$ 时间内先向右做匀减速直线运动，$t_{0} \sim 2t_{0}$ 时间内向右做匀加速直线运动。加速度 $a$ 的大小不变，方向始终向左。

由牛顿第二定律可得：

$F - F_{\text{安}} = -ma$

$F = -ma + F_{\text{安}}$

可见，$F-t$ 图像的斜率与 $F_{\text{安}}$ 的斜率大小相同，并在 $F_{\text{安}}$ 大小图线的基础上向下平移 $ma$ 的距离。

正确选项为 C。
}

\item
%% number: 13
%% paperName: 2023年上海高考第 13 题 4 分
%% typeId: 3
%% type: 填空题
%% chapter: 气体性质
%% point: 分子动理论
%% method: 无
%% score: 4
%% degree: 501
%% duplicateId: 0
%% body:
装在绝热密闭容器中的气体随容器做直线运动，突然停止时其直线运动的动能变为零，则气体温度将\tkanswer[1.2cm]{升高}，气体分子与器壁碰撞的剧烈程度\tkanswer[1.2cm]{增加}。
%% answer: 升高，增加
%% memo:
\memoanswer{
本题原卷及材料中未提供详解，待补充。
}

\item
%% number: 14
%% paperName: 2023年上海高考第 14 题 4 分
%% typeId: 3
%% type: 填空题
%% chapter: 曲线运动
%% point: 天体运动
%% method: 无
%% score: 4
%% degree: 751
%% duplicateId: 0
%% body:
月球在地球引力作用下绕地球做匀速圆周运动，周期为 $T$。月地距离为 $r$，月球质量为 $m$。月球绕地球运动的线速度大小为\tkanswer[2.5cm]{$\frac{2\pi r}{T}$}，地球对月球的万有引力大小为\tkanswer[3cm]{$\frac{4\pi^{2}mr}{T^{2}}$}。
%% answer: $\frac{2\pi r}{T}$，$\frac{4\pi^{2}mr}{T^{2}}$
%% memo:
\memoanswer{
本题原卷及材料中未提供详解，待补充。
}

\item
%% number: 15
%% paperName: 2023年上海高考第 15 题 4 分
%% typeId: 3
%% type: 填空题
%% chapter: 光学
%% point: 光子说
%% method: 无
%% score: 4
%% degree: 870
%% duplicateId: 0
%% body:
核钟是下一代超高精度时间测量工具，通过衰变反应 $\ce{^{238}_{92}U -> ^{234}_{90}Th} +$ \tkanswer[1.5cm]{\ce{^{4}_{2}He}} 可得到核钟所用的 $\ce{^{234}_{90}Th}$。驱动核钟可用光子能量为 $8 \UeV$ 的光波，该光波波长为\tkanswer[2.5cm]{$1.55\times10^{-7}$} $\Um$。（普朗克常数 $h = 6.63\times10^{-34} \UJcdots$，$1 \UeV = 1.6\times10^{-19} \UJ$）
%% answer: $^{4}_{2}$He，1.55$\times$10$^{-7}$
%% memo:
\memoanswer{
由 $E = h\nu = \frac{hc}{\lambda}$ 可得 $\lambda = \frac{hc}{E} = \frac{6.63\times10^{-34}\times 3\times10^{8}}{8\times 1.6\times10^{-19}} \Um \approx 1.55\times10^{-7} \Um$
}

\item
%% number: 16
%% paperName: 2023年上海高考第 16 题 4 分
%% typeId: 3
%% type: 填空题
%% chapter: 光学
%% point: 光子说
%% method: 无
%% score: 4
%% degree: 770
%% duplicateId: 0
%% body:
中国科学家利用光梳技术获得频率为 $\nu_{1}$ 和 $\nu_{2}$（$\nu_{1} < \nu_{2}$）的两束高品质单色光，它们在真空中通过相同距离到达同一双缝的时间\tkanswer[1.2cm]{A}（选填：A．“相等”   B．“不相等”）。这两束单色光通过该双缝后在同一光屏上分别产生干涉条纹，频率为\tkanswer[1.5cm]{$\nu_{1}$}的光产生的条纹间距较大。
%% answer: A，$\nu_{1}$
%% memo:
\memoanswer{
本题原卷及材料中未提供详解，待补充。
}

\item
%% number: 17
%% paperName: 2023年上海高考第 17 题 4 分
%% typeId: 3
%% type: 填空题
%% chapter: 气体性质
%% point: 玻意耳定律
%% method: 无
%% score: 4
%% degree: 320
%% duplicateId: 0
%% body:
为测量大气压强，设计如下实验：一导热气缸开口向上竖直放置，用横截面积为 $S$ 的轻质活塞将一定质量的理想气体封闭在气缸中。在活塞上逐次增加砝码，并记录多组平衡时距气缸底部高度 $h$ 和活塞上砝码总质量 $m$，以 $m$ 为纵轴，$1/h$ 为横轴作图，得到的拟合直线斜率为 $k$，截距为 $b$。据此可得大气压强为\tkanswer[2.5cm]{$-\frac{bg}{S}$}，活塞上无砝码时活塞距气缸底部的高度为\tkanswer[2.5cm]{$-\frac{k}{b}$}。（不计摩擦，重力加速度为 $g$）
\onepicture[w=3cm, align=right]{\includesvg{figs/17}}
%% answer: $-\frac{bg}{S}$，$-\frac{k}{b}$
%% memo:
\memoanswer{
设活塞上无砝码时，活塞距气缸底部高度为 $h_{0}$，大气压强为 $p_{0}$，由玻意耳定律得：

$p_{0}h_{0}S = (p_{0} + \frac{mg}{S})hS$

解得：$m = \frac{p_{0}h_{0}S}{g}\cdot\frac{1}{h} - \frac{p_{0}S}{g}$

$k = \frac{p_{0}h_{0}S}{g}$，$b = -\frac{p_{0}S}{g}$

即 $p_{0} = -\frac{bg}{S}$，$h_{0} = -\frac{k}{b}$
}

\item
%% number: 18
%% paperName: 2023年上海高考第 18 题 10 分
%% typeId: 4
%% type: 实验题
%% chapter: 运动和力
%% point: 验证牛顿第二定律
%% method: 无
%% score: 10
%% degree: 560
%% duplicateId: 0
%% body:
为了研究加速度 $a$ 与 $F$ 的关系，某同学设计的实验装置如图所示。在附有标尺的轨道上固定一光电门传感器，小车上装有一挡光片。测量钩码所受重力的大小，并将其视为小车所受拉力大小。由静止释放小车，记录挡光片经过光电门的挡光时间。改变悬挂的钩码个数，从同一位置释放小车获取多组数据。
\begin{enumerate}
	\item
	该实验\tkanswer[1.2cm]{B}测量小车质量（选择：A．“必须”   B．“不必”）。
	\item
	为了得到 $a-F$ 图像，实验中还需要测量的物理量是：\tkanswer[2.5cm]{挡光片宽度}、\tkanswer[3cm]{光电门到小车的距离}。
	\item
	实际上小车所受拉力略小于钩码所受重力，其原因\tkanswer[8cm]{小车运动时钩码向下做加速运动，由牛顿第二定律可知绳的拉力小于钩码的总重力}。
	\item
	已知细线与轨道平行时小车的加速度为 $a_{1}$，若滑轮位置偏低导致细线与轨道不平行，其他条件不变，小车的加速度 $a_{2}$，则 $a_{1}$\tkanswer[1.2cm]{B}$a_{2}$。（选择：A．“<”       B．“>”）
\end{enumerate}
\onepicture[w=6.5cm, align=right]{\includesvg{figs/18}}
%% answer: 
\jdanswer{
\begin{enumerate}
	\item
	B

	\item
	挡光片宽度，光电门到小车的距离

	\item
	小车运动时钩码向下做加速运动，由牛顿第二定律可知绳的拉力小于钩码的总重力

	\item
	B

\end{enumerate}
}
%% memo:
\memoanswer{
本题原卷及材料中未提供详解，待补充。
}

\item
%% number: 19
%% paperName: 2023年上海高考第 19 题 15 分
%% typeId: 6
%% type: 计算题
%% chapter: 电磁感应
%% point: 线圈过有界磁场
%% method: 无
%% score: 15
%% degree: 601
%% duplicateId: 0
%% body:
如图（a），斜面倾角 $\theta = 30^{\circ}$，斜面所在空间有一边界与斜面底边平行、宽度 $D = 0.40 \Um$、磁感应强度大小 $B = 0.5 \UT$ 的匀强磁场区域，磁场方向垂直斜面向上。单匝矩形闭合金属线框 cdef 质量 $m = 0.10 \Ukg$、总电阻 $R = 0.25 \UO$，cf 边的长度等于磁场宽度。线框在平行斜面向上、垂直于 ef 边的恒定拉力作用下，从斜面底端由静止开始运动，当线框的 cd 边离开磁场区域时撤去拉力。线框运动过程中速度 $v$ 与时间 $t$ 的关系如图（b）所示。已知线框与斜面间的动摩擦因数 $\mu = \frac{\sqrt{3}}{3}$，$g$ 取 $9.8 \Umsq$。求
\begin{enumerate}
	\item
	线框受到的拉力大小 $F$；
	\item
	线框 cd 边的长度 $L$；
	\item
	线框中产生的焦耳热 $Q$。
\end{enumerate}
\twopicture[h=5cm, align=right]
{figs/19a.png}
{figs/19b.png}
%% answer: 
\jdanswer{
\begin{enumerate}
	\item
	$F = 1.48 \UN$

	\item
	$L = 0.5 \Um$

	\item
	$Q = 0.4 \UJ$

\end{enumerate}
}
%% memo:
\memoanswer{
（1）$0 \sim 0.4 \Us$ 内线框做初速为零的匀加速直线运动

$a = \frac{v}{t} = \frac{2}{0.4} \Umsq = 5 \Umsq$

由牛顿第二定律可得

$F - mg\sin 30^{\circ} - \mu mg\cos 30^{\circ} = ma$

$F = mg\sin 30^{\circ} + \mu mg\cos 30^{\circ} + ma$

$= 0.1\times 9.8\times 0.5 + \frac{\sqrt{3}}{3}\times 0.1\times 9.8\times\frac{\sqrt{3}}{2} + 0.1\times 5$

$= 1.48 \UN$

（2）$0.4 \Us$ 后匀速，有

$F - mg\sin 30^{\circ} - \mu mg\cos 30^{\circ} - F_{\text{安}} = 0$

$F - mg\sin 30^{\circ} - \mu mg\cos 30^{\circ} - BIL = 0$

$F - mg\sin 30^{\circ} - \mu mg\cos 30^{\circ} - \frac{B^{2}L^{2}v}{R} = 0$

$\frac{B^{2}L^{2}v}{R} = 0.5$，解得 $L = 0.5 \Um$

（3）线框只有在穿过磁场区域的过程中才会产生热量，其电流强度

$I = \frac{BLv}{R} = \frac{0.5 \times 0.5 \times 2}{0.25} \UA = 2 \UA$

用时

$t = \frac{2D}{v} = \frac{2 \times 0.4}{2} \Us = 0.4 \Us$

产生的焦耳热

$Q = I^{2}Rt = 2^{2}\times 0.25\times 0.4 \UJ = 0.4 \UJ$

（使用 $Q = W_{\text{安克}}$ 亦可）

撤去外力后线框做匀减速直线运动，当速度降为零时，由于

$mg\sin 30^{\circ} = \mu mg\cos 30^{\circ}$

所以不会下滑。
}

\item
%% number: 20
%% paperName: 2023年上海高考第 20 题 15 分
%% typeId: 6
%% type: 计算题
%% chapter: 运动和力
%% point: 动量和能量
%% method: 无
%% score: 15
%% degree: 493
%% duplicateId: 0
%% body:
如图，一质量 $m_{1} = 0.15 \Ukg$ 的小球 P 由长 $l = 1.2 \Um$ 的轻质细线悬挂于 O 点。一质量 $m_{2} = 0.1 \Ukg$ 的小滑块 Q 静置于水平轨道 A 处，位于 O 点正下方。将 P 向左拉开由静止释放，P 在与 Q 碰撞前瞬间的向心加速度 $a = 1.6 \Umsq$，碰撞后瞬间 P 的速度是碰撞前瞬间速度的 $1/5$。滑块与轨道间动摩擦因数 $\mu = 0.28$，碰撞前后小球与滑块的总动能不变，$g$ 取 $9.8 \Umsq$。求：
\begin{enumerate}
	\item
	P 与 Q 碰撞后瞬间物块 Q 的速度大小 $v$；
	\item
	P 与 Q 碰撞后第一次回到最低点时，Q 与 A 点之间的距离 $s$。
\end{enumerate}
\onepicture[w=5cm, align=right]{\includesvg{figs/20}}
%% answer: 
\jdanswer{
\begin{enumerate}
	\item
	$v = 1.663 \Ums$

	\item
	$s = 0.504 \Um$

\end{enumerate}
}
%% memo:
\memoanswer{
（1）由圆周运动向心加速度

$a = \frac{v_{P}^{2}}{l}$

解得 $v_{P} = \frac{4\sqrt{3}}{5} \Ums$

碰撞后小球 P 的速度变为

$v_{P}' = \frac{1}{5}v_{P} = \frac{4\sqrt{3}}{25} \Ums$

碰撞过程系统机械能守恒，有

$\frac{1}{2}m_{1}v_{P}^{2} = \frac{1}{2}m_{1}v_{P}'^{2} + \frac{1}{2}m_{2}v_{Q}^{2}$

解得

$v_{Q} = \frac{24\sqrt{3}}{25} \Ums \approx 1.663 \Ums$

（2）设 P 碰撞后摆动到最高点的高度为 $h$，由机械能守恒可得

$\frac{1}{2}m_{1}v_{P}'^{2} = m_{1}gh$

$h \approx 0.004 \Um$

对应的最大摆角

$\theta = \arccos\frac{l - h}{l} \approx 4.68^{\circ} < 5^{\circ}$

故 P 的摆动可看作简谐振动，P 与 Q 碰撞后再次回到 A 点的时间为

$t_{P} = \frac{T}{2} = \pi\sqrt{\frac{l}{g}} \approx 1.10 \Us$

对 Q 物体，根据牛顿第二定律，有

$f = \mu m_{2}g = m_{2}a$

可得

$a = \mu g = 2.744 \Umsq$

Q 的运动时间

$t_{Q} = \frac{v_{Q}}{a} = \frac{1.663}{2.744} \Us \approx 0.606 \Us < t_{P}$

故 Q 在上述过程中已经停下，其位移为

$s = \frac{v_{Q}^{2}}{2a} = 0.504 \Um$
}

\end{enumerate}

\end{document}
